@@M0@@ on the category of graded @@M1@@-Modules is compatible with tensor product; the same remains true of the functor obtained on the derived categories.

\sgaStableTarget{sga6.e06.sec.2.4}2.4. For every integer @@M2@@, and every @@M3@@-Module @@M4@@, put:
\sgaStableTarget{sga6.e06.eq.2-4-1}@@M5@@
The functor @@M6@@ thus defined is left exact. Since the category of @@M7@@-Modules has enough injectives, we can define a derived functor:
\sgaStableTarget{sga6.e06.eq.2-4-2}@@M8@@

If @@M9@@ is a locally free @@M10@@-Module of finite type, there exists an integer @@M11@@ such that, for @@M12@@, @@M13@@ is @@M14@@-acyclic. To see this, it suffices to show that, for @@M15@@ sufficiently large, @@M16@@ is acyclic for the functors @@M17@@. Since tensoring with @@M18@@ is an exact functor, this follows immediately from the spectral sequence of composite functors and from 1.3.

2.5. We propose to construct, for every @@M19@@-Module @@M20@@, a homomorphism functorial in @@M21@@:
\sgaStableTarget{sga6.e06.eq.2-5-1}@@M22@@
For this, we define the restriction of @@M23@@ to the @@M24@@ (2.3.2). By construction we have, on putting @@M25@@:
@@M26@@
To define a homomorphism:
@@M27@@
